\section{语料与口径}

\subsection{语料与抽取产物}

本文语料是通俗明史叙事《明朝那些事儿》全七部（2006 年起连载出版），
以及在其之上抽取的结构化知识图谱。图谱不是原书本身，而是原书的一次投影；
本篇的四个量全部建立在这一投影之上，第 6 节会反复回到这一点。

规模如表 \ref{tab:corpus} 所示。人物 1231、事件 1106、关系 2324、地点 563。
原书正文不外发，本章不做全文检索；事件类别的抽取由前序工序完成，
本篇只使用其分类结果，不重跑抽取，也不改写任何一条类别标签。

\begin{table}[!t]
\centering
\caption{语料与结构化抽取产物规模}
\label{tab:corpus}
\small
\begin{tabular}{lr}
\toprule
项目 & 数量 \\
\midrule
人物 & 1231 \\
事件 & 1106 \\
关系边 & 2324 \\
地点 & 563 \\
\bottomrule
\end{tabular}
\end{table}

\subsection{三条数据纪律}

本篇沿用项目统一的三条纪律，它们直接决定第 5 节的点数正确性。

\begin{itemize}
\item \textbf{裸名，无 id。} 关系记录的 from / to 存的是裸人物名（如「朱棣」），
人物表里没有 id 字段。凡按名字查表的地方一律用裸名，禁止按 id 拼串；
否则 2324 条边里会有成批的尸体边（端点落在人物表外）。
本篇剔除端点缺失或两端同名的记录后，带边顶点为 741 个（第 5.3 节）。
\item \textbf{坐标不按假值判空。} 地理维度在本篇不使用，因此这条不参与计算，
但凡涉及经纬度的地方必须写 `is None' 判空，不能写 `if not lat'——
0 是合法经纬度，写成假值会把坐标原点上的地点整批丢掉（这是本项目踩过的坑）。
\item \textbf{截断桶必须改名。} 任何 min(v, N) 的截断桶不得沿用被截断的原键名，
必须改名或另立一个 geN 汇总桶。历史上正是这一步把「degree=10 → 143 人」
这类假象塞进过统计；本篇的度数直方图把 13 及以上合成一个桶（记为 $\geq$13），
桶名在表 \ref{tab:degree} 里单独一行，不与 0--12 的十三个桶共名。
\end{itemize}

\subsection{时间口径}

事件年份只取 1300—1700 区间内可解析为整数的记录，得 1006 件，
其余 100 件无年份或不可解析，本篇一律不按年份对待，也不默认补零。
这 1006 件落在 1328—1645 之间的 218 个不同年份上，
其余年份在该口径下为 0（见图 \ref{fig:ts}）。
时间轴的起点在明朝立国（1368）之前，这是原书叙事自身的跨度，
不是本篇的切分错误；第 5.5 节的三个时期均落在 1368 之后。
